\documentclass[9pt,conference]{IEEEtran}
\usepackage[utf8]{inputenc}
\usepackage[spanish,english]{babel}
\usepackage{graphicx}
\usepackage{booktabs}
\usepackage{caption}
\usepackage{subcaption}
\usepackage{url}
\usepackage{hyperref}
\usepackage{tabularx}
\usepackage{booktabs}
\usepackage{threeparttable}

\title{Food Waste Reduction Platform: Robust System Design and Project Management}

\author{
Luna Alejandra Sandoval Rodríguez \\
\textit{Project Lead \& Systems Analyst} \\
Universidad Distrital Francisco José de Caldas \\
Bogotá, Colombia \\
20241020053
\and
Cristian Camilo Bonilla Lizarazo \\
\textit{Frontend Developer \& Systems Analyst} \\
Universidad Distrital Francisco José de Caldas \\
Bogotá, Colombia \\
20241020015
\and
Nicolás Rodríguez Granados \\
\textit{Backend Developer \& Systems Analyst} \\
Universidad Distrital Francisco José de Caldas \\
Bogotá, Colombia \\
20241020037
\and
Juan Sebastián Bravo Rojas \\
\textit{Database Specialist \& Systems Analyst} \\
Universidad Distrital Francisco José de Caldas \\
Bogotá, Colombia \\
20241020004
}

\begin{document}

\maketitle

\section{Project Management Documentation}

\subsection{Project Charter}

\textbf{Project Title:}

Food Waste Reduction Platform: A Real-Time Geospatial Redistribution System.

\textbf{Problem Statement and Justification:}

Despite Colombia generating 9.7 million tons of food waste annually, local coordination gaps prevent the efficient redistribution of high-quality surplus within short freshness windows. This project bridges that gap by implementing a high-concurrency matching engine that prioritizes social equity and geographical proximity, directly supporting the circular economy objectives of the Bogotá university community.


\subsubsection{SMART Objectives}
\begin{itemize}
    \item \textbf{Specific:} To engineer and validate a mobile ecosystem connecting donor nodes (cafeterias/restaurants) with recipient nodes (students/charities) within a strictly enforced 3 km radius of the Carrera 8 \#40-62 headquarters.
    \item \textbf{Measurable:} To achieve a Surplus Recovery Rate (SRR) of $\ge 80\%$, a backend algorithmic latency of $< 2$ seconds, and a Verified Pickup Rate (VPR) of $\ge 85\%$ under simulated load.
    \item \textbf{Achievable:} Leveraging the validated Workshop 2 stack: NestJS (logic), PostGIS (spatial), Redis (concurrency), and React Native (interface).
    \item \textbf{Relevant:} Directly addresses the 17 design requirements and 10 critical risks identified in the Workshop 1 systems diagnostic.
    \item \textbf{Time-bound:} Final system validation and performance benchmarking to be concluded by June 12, 2026.
\end{itemize}

\subsubsection{Comprehensive Team Structure and Accountability (RACI)}


To ensure project success, the team operates under a RACI Matrix (Responsible, Accountable, Consulted, Informed), ensuring no ambiguity in deliverable ownership.

See Table I: Accountability Framework (RACI)

\begin{table}[h]
\centering
\small
\begin{threeparttable}
    \caption{Accountability Framework (RACI)}
    \begin{tabularx}{\columnwidth}{@{} X cccc @{}}
        \toprule
        \textbf{Activity} & \textbf{Luna} & \textbf{Nico} & \textbf{Cris} & \textbf{Juan} \\ \midrule
        Project Lead \& Integration & A/R & I & I & I \\
        Backend \& Matching Engine  & I & A/R & C & C  \\ 
        Frontend \& UX Design       & I & C & A/R & I \\ 
        Quality Assurance (QA)      & R & R  & R  & R   \\
        \bottomrule
    \end{tabularx}
    \begin{tablenotes}
        \small
        \item \textbf{Legend:} R: Responsible; A: Accountable; C: Consulted; I: Informed.
    \end{tablenotes}
\end{threeparttable}
\end{table}

\subsubsection{Project Timeline and Critical Path}

The project follows a 16-week cycle divided into four major iterations. The Critical Path currently lies in the integration of the matching engine with real-time push notifications.
\begin{itemize}
    \item \textbf{M1: System Diagnostic (Weeks 1-4):} Problem definition, IPO modeling, and boundary setting. [COMPLETED]
    \item \textbf{M2: Architectural Design (Weeks 5-8):} Component diagrams, technology stack selection, and risk mitigation planning. [COMPLETED]
    \item \textbf{M3: Development \& Optimization (Weeks 9-12):} Sprint-based coding, PostGIS implementation, and reliability scoring logic. [IN PROGRESS]
    \item \textbf{M4: Validation \& Final Report (Weeks 13-16):} Stress testing, experimental validation, and presentation of the "Best-in-Class" model.
\end{itemize}

\subsubsection{Resource Planning and Budgetary Considerations}

As an academic initiative, the budget is calculated in Human-Hours (HH) and Cloud Resource Credits:
\begin{itemize}
    \item \textbf{Human Capital:} 4 Systems Engineers x 10 hours/week x 16 weeks = 640 total HH.
    \item \textbf{Technological Assets:} Expo/React Native (Open Source), PostgreSQL/PostGIS (Open Source), and Firebase (Free Tier).
    \item \textbf{Infrastructure Cost (Theoretical):} Estimated \$50 USD/month for production-grade hosting on AWS/Google Cloud for a pilot-scale deployment.
\end{itemize}

\subsubsection{Stakeholder Communication Strategy}To maintain alignment with both internal and external stakeholders, the following communication matrix is enforced:
\begin{itemize}
    \item \textbf{Technical Syncs (Internal):} Bi-weekly agile stand-ups to review GitHub issues and code reviews.
    \item \textbf{Academic Oversight:} Monthly progress reports submitted via the course platform to Eng. Carlos Andrés Sierra.
    \item \textbf{External Validation:} Semi-structured interviews with 3 local cafeteria managers to validate the "Same-day 2-4h window" feasibility.
\end{itemize}

\subsubsection{Scope and Constraints}
\textbf{Functional Boundary:} The system is restricted to digital coordination. It does not manage cold-chain logistics, nor does it assume legal liability for food quality, which remains the donor's responsibility as per the system's "Terms of Use" architecture.
\textbf{Assumptions:}
\begin{enumerate}
    \item Continuous availability of GPS services for both donor and recipient.
    \item Stable internet connectivity within the 3 km university radius.
    \item Users (students/collectors) possess smartphones compatible with the Expo Go environment.
\end{enumerate}


\begin{thebibliography}{10}
\bibitem{w1} Workshop 1 – Food Waste Reduction Mobile Application, A Rigorous System Analysis, Universidad Distrital Francisco José de Caldas, 2026.
\bibitem{w2} Workshop 2 – System Design Document, Universidad Distrital Francisco José de Caldas, 2026.
\bibitem{ieee} IEEE Standards Association, ``IEEE 730-2014 - IEEE Standard for Software Quality Assurance Processes,'' 2014.
\end{thebibliography}

\end{document}
